\chapter{Refusing the Postulate}\label{app:refusing}

The book makes one explicit interpretive commitment: the postulate of Chapter 6, that no cut is privileged. Chapter 7 said that a reader could refuse it and keep the mathematical results and universal construction. This appendix says plainly what that refusal changes, and what it does not decide.

Start with what survives. Bayes' theorem is unchanged. Cox-style representation results retain their assumptions. Kraft still constrains prefix-code lengths. Fair program bits and a reference universal machine still define $M$, and dominance, convergence, and invariance retain their scoped statements. None mentions the postulate in its proof. Most of Chapter 2's inversion survives as well: generating the Library is still cheap, locating anything in it is still expensive, and complexity is still uncomputable. You may use $M$ as a theoretical predictor without giving it an ontological reading. You may also prefer another coherent prior for a finite problem; Part I never proved that $M$ is uniquely mandatory.

What you give up is the postulate's refusal to privilege any decomposition, and what you take on is at least the possibility, and usually the assertion, that some carving is the structure's own. The book declined to assert either way, since nothing inside could certify a joint or its absence (the machine-relativity of Chapter 5 for the one direction, and Chapter 2's one-way certificate, stated formally in Appendix~\ref{app:bedrock}, for the other), and a reader who refuses it and asserts a joint keeps every theorem, the construction, and all three of the principles the body names: no description contains its center, weights come after cuts, and a found name certifies itself while an absence cannot be certified. Refusal does not force the opposite package. You may be a one-world realist, a structural realist who accepts some natural decompositions, an instrumentalist about $M$, or something else. Likewise, treating the self as a useful local model can be motivated by cognitive science and Parfitian arguments without any multiverse. The fork is not ``accept the postulate and all worlds'' versus ``refuse it and one ordinary self.'' Several independent questions remain.

The book does not argue you out of those alternatives. Chapter 6 compared the description length of selecting one candidate with enumerating program space, while warning that code length is not automatically ontological parsimony. Keeping an ordinary metaphysics is not a failure of nerve or a misreading. If you have read this far and still prefer the world you started with, you have understood the fork. Keep the theorems, use or decline the construction, and choose the metaphysics separately.
